\documentclass[diss,capa,times,english]{ppgci}
%Remova english no \documentclass se optar pelo português
%Substitua times por arial caso queira mudar a fonte

\usepackage{microtype} %Melhorias de justificação e espaçamento tipográfico
\usepackage{enumerate} %Personalizar listas (a, b, c, i, ii, etc.)
\usepackage{pdfpages} %Inserir páginas de arquivos PDF no documento
\usepackage{multirow} %Mesclar linhas em tabelas
\usepackage{comment}  %Adicionar e ocultar blocos de comentário
\usepackage{longtable} %Tabelas longas com quebra automática de página
\usepackage{booktabs} %Melhorar estilo de tabelas
\usepackage{pdflscape} %Permitir rotacionar a página inteira
\usepackage{rotating} %Permitir rotacionar tabela 
\usepackage{float} %Controlar de posicionamento de floats

%Inserção da bibliografia
\addbibresource{bibliografia.bib} 

\hypersetup{
    hidelinks, %Remove coloração e caixas
    unicode=true, %Permite acentuação no bookmark
    linktoc=all %Habilita link no nome e página do sumário
}

%Informações sobre a Pesquisa
\areapesquisa{Gestão e Economia da Construção}
\areapesquisaeng{Management and Economics of Construction} %Insira a área de pesquisa em inglês

\titulo{Título do Trabalho} 
\tituloeng{Title of the Work} %Insira o título em inglês

\author{Wojciechowski}{Anderson Martins} 
\email{andersoncivil3@gmail.com}

\advisor[Prof.]{Isatto}{Eduardo Luis}{Dr., UFRGS, Brasil}
%\advisor[Prof.]{Isatto}{Eduardo Luis}{Dr., UFRGS, Brasil}
%\coadvisor[Prof.]{Formoso}{Carlos Torres}{Dr., UFRGS, Brasil} %Descomente caso tenha coorientador
%\coadvisor[Prof.]{Formoso}{Carlos Torres}{Dr., UFRGS, Brasil}

%Palavras-chave em PT_BR
%use letras minúsculas seguidas por ;
%a última terminada por .
\keyword{palavrachave-um;}
\keyword{palavrachave-dois;}
\keyword{palavrachave-tres;}
\keyword{palavrachave-quatro.}

%Palavras-chave em EN_US
\keywordeng{keyword-one;}
\keywordeng{keyword-two;}
\keywordeng{keyword-three;}
\keywordeng{keyword-four.}

\begin{document}

%\renewcommand{\advisorname}{Orientadora}           %Descomente caso tenha orientadora
%\renewcommand{\coadvisorname}{Coorientadora}      %Descomente caso tenha coorientadora

\maketitle 

\sloppy

\fichacatalografica

%Composição da Banca Examinadora
\coordenador{Lélio Antônio Teixeira Brito}
\begin{aprovacao}{30}{2}{2026}
\noindent \textbf {Prof. XXX (Universidade)}\\
Título, Universidade, País.\\[1cm] 

\noindent \textbf {Prof. XXX (Universidade)}\\
Título, Universidade, País.\\[1cm] 

\noindent \textbf {Prof. XXX (Universidade)}\\
Título, Universidade, País.
\end{aprovacao}

%Opcional
\begin{dedicatoria}
  Dedico\ldots 
\end{dedicatoria}

%Opcional
\begin{agradecimentos}
  Agradeço\ldots 
\end{agradecimentos}

%Opcional
\begin{epigrafe}
    \epigrafetexto{Começar é a parte mais importante do trabalho.}
    \epigrafeautor{Platão}
\end{epigrafe}

%Opcional
%\forcabrabstract %Caso o trabalho seja apresentado em inglês, a seção “RESUMO” é dispensável. Se quiser mostrar o resumo em português também, descomente.

\hyphenpenalty=10000 %Justificar um texto sem utilizar hifenização
\exhyphenpenalty=10000 %Controla quebras em hífens explícitos no texto

%Resumo em Português (no máximo 500 palavras)
\ifshowabstractpt
\begin{abstract}
Bla blabla blablabla bla.  Bla blabla blablabla bla.  Bla blabla
blablabla bla.  Bla blabla blablabla bla.  Bla blabla blablabla bla.
Bla blabla blablabla bla.  Bla blabla blablabla bla.  Bla blabla
blablabla bla.  Bla blabla blablabla bla.  Bla blabla blablabla bla.
Bla blabla blablabla bla.  Bla blabla blablabla bla.  Bla blabla
blablabla bla.  Bla blabla blablabla bla.  Bla blabla blablabla bla.
Bla blabla blablabla bla.  Bla blabla blablabla bla.  Bla blabla
blablabla bla.  Bla blabla blablabla bla.  Bla blabla blablabla bla.
Bla blabla blablabla bla.
\end{abstract}
\fi

%Resumo em Inglês (no máximo 500 palavras)
\begin{englishabstract}
Bla blabla blablabla bla.  Bla blabla blablabla bla.  Bla blabla
blablabla bla.  Bla blabla blablabla bla.  Bla blabla blablabla bla.
Bla blabla blablabla bla.  Bla blabla blablabla bla.  Bla blabla
blablabla bla.  Bla blabla blablabla bla.  Bla blabla blablabla bla.
Bla blabla blablabla bla.  Bla blabla blablabla bla.  Bla blabla
blablabla bla.  Bla blabla blablabla bla.  Bla blabla blablabla bla.
Bla blabla blablabla bla.  Bla blabla blablabla bla.  Bla blabla
blablabla bla.  Bla blabla blablabla bla.  Bla blabla blablabla bla.
Bla blabla blablabla bla.
\end{englishabstract}

%Lista de Figuras
\listoffigures

%Lista de Tabelas
\listoftables

%Lista de Quadros - AMW/Marilton
\listofquadros %Se optar pelo inglês não existirá lista de quadros, pois documentos em inglês não há distinção entre tabelas e quadros

%lista de abreviaturas e siglas
\begin{listofabbrv}{UFRGS}%coloque aqui a maior sigla para ajustar a distância
    \item[ABNT] Associação Brasileira de Normas Técnicas
    \item[DSR] Design Science Research
    \item[EE] Escola de Engenharia
    \item[NORIE] Núcleo Orientado para a Inovação da Construção
    \item[PPGCI] Programa de Pós-Graduação em Engenharia Civil: Construção e Infraestrutura
    \item[UFRGS] Universidade Federal do Rio Grande do Sul
\end{listofabbrv}

%Sumário
\tableofcontents

\chapter{Introdução} %Inglês: Introduction
Este capítulo apresenta o contexto e a justificativa do tema proposto, os pressupostos da pesquisa, o problema de pesquisa (lacuna do conhecimento), as questões de pesquisa e os objetivos da pesquisa.

Alguns comandos de formatação tipográfica que podem ser utilizados ao longo do texto. No entanto, em textos acadêmicos formais, recomenda-se o uso moderado desses recursos, para manter a padronização visual do documento. Entre os principais comandos estão: \emph{texto em itálico}, \textbf{texto em negrito}, \textbf{\emph{texto em itálico e negrito}}, \underline{texto sublinhado}, \textsc{texto em versalete}, \texttt{texto monoespaçado} e \small{texto em tamanho menor}.

Também são apresentados exemplos de procedimentos para realizar citações no texto. Quando o autor é citado ao final da frase, a referência deve ser apresentada entre parênteses, indicando o sobrenome e o ano de publicação, separados por vírgula~\cite{Sacks2018}. Quando o nome do autor é mencionado no corpo do texto, utiliza-se o formato narrativo~\textcite{Ensafi2021}. 

Para citações diretas, deve-se indicar a página consultada~\cite[p.45]{Borrmann2018}. No caso de múltiplas obras de um mesmo autor publicadas no mesmo ano, devem ser adotados identificadores alfabéticos~\cite{ABNT2022a,ABNT2022b}. Já para múltiplas obras de um mesmo autor publicadas em anos diferentes~\cite{ABNT2022a,ABNT2024}.

Para citações de diferentes autores em uma mesma referência, elas devem ser apresentadas conjuntamente~\cite{Isatto2005,Ensafi2021,buildingSMART2024}. Por fim, para citações indiretas de outra fonte (citação da citação), utiliza-se o comando~\apud{Wojciechowski2025}{Wojciechowski2024}.

\section{Contexto de Pesquisa} %Inglês: Research Context

Bla blabla blablabla bla.  Bla blabla blablabla bla.  Bla blabla
blablabla bla. Bla blabla blablabla bla.  Bla blabla blablabla bla.  Bla blabla
blablabla bla.  Bla blabla
blablabla bla.  Bla blabla blablabla bla.  Bla blabla blablabla bla. 
Bla blabla blablabla bla.  Bla blabla blablabla bla.  Bla blabla
blablabla bla.  Bla blabla blablabla bla.  Bla blabla blablabla bla.
Bla blabla blablabla bla.  Bla blabla blablabla bla.  Bla blabla
blablabla bla.  Bla blabla blablabla bla.  Bla blabla blablabla bla.
Bla blabla blablabla bla.  Bla blabla blablabla bla.  Bla blabla
blablabla bla.  Bla blabla blablabla bla.  Bla blabla blablabla bla
Conforme a Figura~\ref{Figuras/castelinhonorie}.

% *** FIGURA ***
\begin{figure}[!htb]
    \centering 
    \caption{Nome da figura}
    \includegraphics[angle=0,scale=.16]{Figuras/castelinhonorie} \\
    {\small Fonte: elaborada pelo autor.}
    \label{Figuras/castelinhonorie}
\end{figure}
% *** FIGURA ***

No LaTeX, [!htbp] define as opções de posicionamento de elementos flutuantes, como tabelas, quadros e figuras. Cada letra representa uma possibilidade conforme apresentado no Quadro~\ref{qua_posicionamento}.

% *** QUADRO ***
\begin{quadro}[!htbp]
    \centering
    \caption{Opções de posicionamento de elementos flutuantes no \LaTeX.}
    \label{qua_posicionamento}
    \begin{tabular}{|p{3cm}|p{3cm}|p{8cm}|}
        \hline
        \small \textbf{Código} & \small \textbf{Significado} & \small \textbf{Descrição}\\
        \hline
        {\small \texttt{h}} & {\small Here} & {\small Tentar posicionar no local onde o elemento foi inserido no texto.}\\
        \hline
        {\small \texttt{t}} & {\small Top} & {\small Posicionar no topo da página.}\\
        \hline
        {\small \texttt{b}} & {\small Bottom} & {\small Posicionar na parte inferior da página.}\\
        \hline
        {\small \texttt{p}} & {\small Page} & {\small Posicionar em uma página exclusiva para elementos flutuantes.}\\
        \hline
        {\small \texttt{!}} & {\small Override} & {\small Relaxar algumas restrições internas do \LaTeX{} para permitir o posicionamento.}\\
        \hline
        {\small \texttt{H}} & {\small Here (Forçado)} & {\small Fixar obrigatoriamente o elemento no local exato do texto.}\\
        \hline
    \end{tabular}
    \vspace{2pt}
    {\small Fonte: elaborado pelo autor.}
\end{quadro}
% *** QUADRO ***

Bla blabla blablabla bla.  Bla blabla blablabla bla.  Bla blabla
blablabla bla.  Bla blabla blablabla bla.  Bla blabla blablabla bla.
Bla blabla blablabla bla. 
Bla blabla blablabla bla.  Bla blabla blablabla bla.  Bla blabla
blablabla bla.  Bla blabla blablabla bla.  Bla blabla blablabla bla.
Bla blabla blablabla bla.  Bla blabla blablabla bla.  Bla blabla
blablabla bla.  Bla blabla blablabla bla.  Bla blabla blablabla bla.
Bla blabla blablabla bla.  Bla blabla blablabla bla.  Bla blabla
blablabla bla.  Bla blabla blablabla bla.  Bla blabla blablabla bla
Bla blabla blablabla bla.  Bla blabla blablabla bla.  Bla blabla
blablabla bla.  Bla blabla blablabla bla.  Bla blabla blablabla bla.
Bla blabla blablabla bla. 

\section{Problema Prático} %Inglês: Practical Problem

Bla blabla blablabla bla.  Bla blabla blablabla bla.  Bla blabla
blablabla bla.  Bla blabla blablabla bla.  Bla blabla blablabla bla.
Bla blabla blablabla bla.  Bla blabla blablabla bla.  Bla blabla
blablabla bla.  Bla blabla blablabla bla.  Bla blabla blablabla bla.
Bla blabla blablabla bla.  Bla blabla blablabla bla.  Bla blabla
blablabla bla.  Bla blabla blablabla bla.  Bla blabla blablabla bla

\section{Problema de Pesquisa} %Inglês: Research Problem

Bla blabla blablabla bla.  Bla blabla blablabla bla.  Bla blabla
blablabla bla.  Bla blabla blablabla bla.  Bla blabla blablabla bla.
Bla blabla blablabla bla.  Bla blabla blablabla bla.  Bla blabla
blablabla bla.  Bla blabla blablabla bla.  Bla blabla blablabla bla.
Bla blabla blablabla bla.  Bla blabla blablabla bla.  Bla blabla
blablabla bla.  Bla blabla blablabla bla.  Bla blabla blablabla bla.

\section{Questões de Pesquisa} %Inglês: Research Questions

Com base no problema de pesquisa, foi definida a questão principal de pesquisa:

\textbf{Questão principal?}

Como desdobramento da questão principal, foram definidas as questões secundárias: 

\begin{itemize}
    \item[i.] Questão secundária 1?
    \item[ii.] Questão secundária 2?
    \item[iii.] Questão secundária 3?
\end{itemize}

\section{Objetivos de Pesquisa} %Inglês: Research Objectives

Objetivo geral da pesquisa:

\textbf{Objetivo geral.}

Objetivos específicos da pesquisa:

\begin{itemize}
    \item[i.] Objetivo específico 1.
    \item[ii.] Objetivo específico 2.
    \item[iii.] Objetivo específico 3.
\end{itemize}

\section{Estrutura do Trabalho} 

Bla blabla blablabla bla.  Bla blabla blablabla bla.  Bla blabla
blablabla bla.  Bla blabla blablabla bla.  Bla blabla blablabla bla.

\chapter{Revisão bibliográfica} %Inglês: Literature Review

Escrever sobre o que se trata o capítulo.

\section{Uma seção}

Bla blabla blablabla bla.  Bla blabla blablabla bla.  Bla blabla
blablabla bla.  Bla blabla blablabla bla.  Bla blabla blablabla bla.
Bla blabla blablabla bla.  Bla blabla blablabla bla.  Bla blabla
blablabla bla.  Bla blabla blablabla bla.  Bla blabla blablabla bla.
Bla blabla blablabla bla.  Bla blabla blablabla bla.  Bla blabla
blablabla bla.  Bla blabla blablabla bla.  Bla blabla blablabla bla.
Bla blabla blablabla bla.  Bla blabla blablabla bla.  Bla blabla
blablabla bla.  Bla blabla blablabla bla.  Bla blabla blablabla bla.
Bla blabla blablabla bla.

\subsection{Uma subseção}

Bla blabla blablabla bla.  Bla blabla blablabla bla.  Bla blabla
blablabla bla.  Bla blabla blablabla bla.  Bla blabla blablabla bla.
Bla blabla blablabla bla.  Bla blabla blablabla bla.  Bla blabla
blablabla bla.  Bla blabla blablabla bla.  Bla blabla blablabla bla.
Bla blabla blablabla bla.  Bla blabla blablabla bla.  Bla blabla
blablabla bla.  Bla blabla blablabla bla.  Bla blabla blablabla bla.
Bla blabla blablabla bla.  Bla blabla blablabla bla.  Bla blabla
blablabla bla.  Bla blabla blablabla bla.  Bla blabla blablabla bla.
Bla blabla blablabla bla.

\subsubsection{Uma subsubseção}

Bla blabla blablabla bla.  Bla blabla blablabla bla.  Bla blabla
blablabla bla.  Bla blabla blablabla bla.  Bla blabla blablabla bla.
Bla blabla blablabla bla.  Bla blabla blablabla bla.  Bla blabla
blablabla bla.  Bla blabla blablabla bla.  Bla blabla blablabla bla.
Bla blabla blablabla bla.  Bla blabla blablabla bla.  Bla blabla
blablabla bla.  Bla blabla blablabla bla.  Bla blabla blablabla bla.
Bla blabla blablabla bla.  Bla blabla blablabla bla.  Bla blabla
blablabla bla.  Bla blabla blablabla bla.  Bla blabla blablabla bla.
Bla blabla blablabla bla Figura~\ref{Figuras/fig_rotacionada}.

% *** FIGURA ROTACIONADA ***
\begin{landscape}
    \begin{figure}[ht]
        \centering 
        \caption{Exemplo de figura em página rotacionada}
        \includegraphics[angle=0,scale=0.6]{Figuras/fig_rotacionada}
        \label{Figuras/fig_rotacionada}\\
        {\small Fonte: elaborada pelo autor.}
    \end{figure}
\end{landscape}
% *** FIGURA ROTACIONADA ***

\chapter{Método de Pesquisa} %Inglês: Research Method

Este capítulo apresenta o método que será seguido para o desenvolvimento desta pesquisa, incluindo a estratégia de pesquisa, o delineamento da pesquisa, as fontes de evidência a serem utilizadas e o cronograma de atividades previsto.

\section{Estratégia de Pesquisa}

Bla blabla blablabla bla.  Bla blabla blablabla bla.  Bla blabla
blablabla bla.  Bla blabla blablabla bla.  Bla blabla blablabla bla.
Bla blabla blablabla bla.  Bla blabla blablabla bla.  Bla blabla
blablabla bla.  Bla blabla blablabla bla.  Bla blabla blablabla bla.
Bla blabla blablabla bla.  Bla blabla blablabla bla.  Bla blabla
blablabla bla.  Bla blabla blablabla bla.  Bla blabla blablabla bla.

\section{Delineamento da Pesquisa}

Bla blabla blablabla bla.  Bla blabla blablabla bla.  Bla blabla
blablabla bla.  Bla blabla blablabla bla.  Bla blabla blablabla bla.
Bla blabla blablabla bla.  Bla blabla blablabla bla.  Bla blabla
blablabla bla.  Bla blabla blablabla bla.  Bla blabla blablabla bla.
Bla blabla blablabla bla.  Bla blabla blablabla bla.  Bla blabla
blablabla bla.  Bla blabla blablabla bla.  Bla blabla blablabla bla.

\section{Fontes de Evidências}

Bla blabla blablabla bla.  Bla blabla blablabla bla.  Bla blabla
blablabla bla.  Bla blabla blablabla bla.  Bla blabla blablabla bla.
Bla blabla blablabla bla.  Bla blabla blablabla bla.  Bla blabla
blablabla bla.  Bla blabla blablabla bla.  Bla blabla blablabla bla.
Bla blabla blablabla bla.  Bla blabla blablabla bla.  Bla blabla
blablabla bla.  Bla blabla blablabla bla.  Bla blabla blablabla bla.

\chapter{Resultados} %Inglês: Results

Bla blabla blablabla bla.  Bla blabla blablabla bla.  Bla blabla
blablabla bla. Bla blabla blablabla bla.  Bla blabla blablabla bla.  Bla blabla
blablabla bla.  Bla blabla
blablabla bla.  Bla blabla blablabla bla.  Bla blabla blablabla bla.
Bla blabla blablabla bla.  Bla blabla blablabla bla.  Bla blabla
blablabla bla.  Bla blabla blablabla bla.  Bla blabla blablabla bla.
Bla blabla blablabla bla. 

Uma tabela apresenta dados quantitativos, numéricos ou estatísticos submetidos a tratamento analítico. A Tabela~\ref{tab_dados} apresenta um exemplo de como criar uma tabela.

% *** TABELA ***
\begin{table}[ht]
    \centering
    \caption{Detalhes do artigo, status, periódicos e fator de impacto.}
    \label{tab_dados}
    \begin{tabular}{x{3,5cm}x{3,5cm}x{3,5cm}x{3,5cm}} %Largura útil = 14,3cm
        \hline
            {\small \textbf{Coluna 1}} & {\small \textbf{Coluna 2}} & {\small \textbf{Coluna 3}} & {\small \textbf{Coluna 4}}\\
            \hline
            {\small Linha 1} & {\small a} & {\small x} & {\small $\alpha$}\\
            {\small Linha 2} & {\small b} & {\small y} & {\small $\beta$}\\
            {\small Linha 3} & {\small c} & {\small z} & {\small $\gamma$}\\
        \hline
    \end{tabular}
    \vskip 2pt
    {\centering \small Fonte: elaborada pelo autor.}   
\end{table}
% *** TABELA ***

Um quadro apresenta dados qualitativos, informações textuais, resumos ou comparações conceituais (sem foco em cálculos numéricos). O Quadro~\ref{qua_informacao} apresenta um exemplo de como criar um quadro. Vale ressaltar que, em documentos em inglês, não há distinção entre tabelas e quadros; deve-se utilizar apenas tabelas. 

% *** QUADRO ***
\begin{quadro}[!htbp]
    \centering
    \caption{Detalhes dos artigos, status, periódicos e fator de impacto.}
    \label{qua_informacao}
    \begin{tabular}{|x{5cm}|x{3cm}|x{3cm}|x{3cm}|}
        \hline
        {\small \textbf{Artigo}} & {\small \textbf{1}} & {\small \textbf{2}} & {\small \textbf{3}}\\
        \hline
        {\small Título do artigo} & {\small Título do Artigo 1} & {\small Título do Artigo 2} & {\small Título do Artigo 3}\\
        \hline
        {\small Status} & {\small Publicado} & {\small Submetido} & {\small Manuscrito}\\
        \hline
        {\small Periódico} & {\small Automation in Construction} & {\small Automation in Construction} & {\small Automation in Construction}\\
        \hline
        {\small Fator de impacto (IF)} & {\small 11,5} & {\small 11,5} & {\small 11,5}\\
        \hline
    \end{tabular}
    \vspace{2pt}
    {\small Fonte: elaborado pelo autor.}
\end{quadro}
% *** QUADRO ***

Mais detalhes dos resultados nos Apêndices~\ref{apen_A},~\ref{apen_B} e no Anexo~\ref{anex_A}.

\chapter{Conclusões (ou Considerações Finais)} %Inglês: Conclusions (or Final Considerations)

Bla blabla blablabla bla.  Bla blabla blablabla bla.  Bla blabla
blablabla bla. Bla blabla blablabla bla.  Bla blabla blablabla bla.  Bla blabla
blablabla bla.  Bla blabla
blablabla bla.  Bla blabla blablabla bla.  Bla blabla blablabla bla.
Bla blabla blablabla bla.  Bla blabla blablabla bla.  Bla blabla
blablabla bla.  Bla blabla blablabla bla.  Bla blabla blablabla bla.
Bla blabla blablabla bla. 

%Imprime as referências bibliográficas 
\printbibliography      % biblatex

%Apêndices (Opcional) - Material produzido pelo autor
\apendices
\chapter{Um Apêndice} \label{apen_A} 

\includepdf[pages=-]{Figuras/apendiceA}

\chapter{Outro Apêndice} \label{apen_B}

Bla blabla blablabla bla.  Bla blabla blablabla bla.  Bla blabla
blablabla bla.  Bla blabla blablabla bla.  Bla blabla blablabla bla.
Bla blabla blablabla bla.  Bla blabla blablabla bla.  Bla blabla
blablabla bla.  Bla blabla blablabla bla.  Bla blabla blablabla bla.
Bla blabla blablabla bla.  Bla blabla blablabla bla.  Bla blabla
blablabla bla.  Bla blabla blablabla bla.  Bla blabla blablabla bla.
Bla blabla blablabla bla.  Bla blabla blablabla bla.  Bla blabla
blablabla bla.  Bla blabla blablabla bla.  Bla blabla blablabla bla.
Bla blabla blablabla bla.

%Anexos (Opcional) - Material produzido por outro
\anexos
\chapter{Forma de dividir o título do Anexo e Apêndice \texorpdfstring{\\}{} em duas linhas} \label{anex_A}

\includepdf[pages=-]{Figuras/anexoA}

\end{document}

